% cumcm-toolkit/templates/paper.tex
% 编译：xelatex paper.tex （跑两趟；ctex 宏包不支持 pdflatex）
\documentclass[12pt]{article}
\usepackage{ctex}
\usepackage{amsmath, amssymb}
\usepackage{graphicx}
\usepackage{booktabs}
\usepackage[colorlinks=true, linkcolor=black]{hyperref}
\usepackage{geometry}
\geometry{a4paper, margin=2.5cm}

\title{论文标题}
\author{}
\date{}

\begin{document}
\maketitle

\begin{abstract}
摘要正文。一段话说清问题、方法、关键结果、亮点。

\vspace{1em}
\noindent\textbf{关键词：} 关键词1；关键词2；关键词3（4-8 个，参见 analysis/writing-style-guide.md）
\end{abstract}

\section{问题重述}

\section{问题分析}

\section{模型假设}

\begin{enumerate}
  \item 假设一
\end{enumerate}

\section{符号说明}

\begin{table}[htbp]
  \centering
  \caption{符号说明}
  \begin{tabular}{cl}
    \toprule
    符号 & 含义 \\
    \midrule
    $Z$ & 目标函数 \\
    $x_i$ & 决策变量 \\
    \bottomrule
  \end{tabular}
\end{table}

\section{模型的建立与求解}

\subsection{问题一}

\section{灵敏度分析}

\section{模型评价与推广}

\subsection{模型的优点}

\subsection{模型的缺点}

\begin{thebibliography}{9}
\bibitem{ref1} 作者. 文献标题[J]. 期刊名, 年份.
\end{thebibliography}

\appendix
\section{支撑文件/材料清单}
% 2022 年起获奖论文标配，见 analysis/judge-deductions.md 对应检查项

\section{程序代码}

\end{document}
